% block2_btcs.tex
%
% Phase 3, Block 2: implicit BTCS scheme + Thomas algorithm.
%
% Self-contained writeup. Compiles standalone with pdflatex.

\documentclass[11pt,a4paper]{article}

\usepackage[T1]{fontenc}
\usepackage[utf8]{inputenc}
\usepackage{amsmath, amssymb, amsthm}
\usepackage{geometry}
\usepackage{hyperref}
\usepackage{enumitem}

\geometry{margin=1in}

\newtheorem{proposition}{Proposition}[section]
\newtheorem{theorem}[proposition]{Theorem}
\newtheorem{definition}[proposition]{Definition}
\newtheorem{remark}[proposition]{Remark}

\title{Phase 3, Block 2:\\The implicit BTCS scheme\\and the Thomas algorithm}
\author{Quant Finance Portfolio}
\date{}

\begin{document}
\maketitle

\section{Setting and motivation}

Block~1 implemented the explicit FTCS scheme, which is conditionally
stable: $\alpha = (\sigma^2/2)\Delta\tau/\Delta x^2 \le 1/2$ forces
$M \gtrsim N^2$ time steps and $O(N^3)$ total work. This block
implements BTCS (Backward-Time, Centred-Space), an
\emph{unconditionally stable} scheme: any $\Delta\tau > 0$ produces a
convergent solution, decoupling $M$ from $N$.

The cost of unconditional stability is that each time step requires
solving a linear system. Block~2 has therefore two deliverables:
\begin{enumerate}[leftmargin=2em]
\item the BTCS pricer for European options;
\item the \emph{Thomas algorithm}, a specialised Gaussian elimination
that solves tridiagonal systems in $O(N)$ work.
\end{enumerate}
Thomas is implemented from scratch as a standalone module
(\texttt{quantlib.thomas} in Python, \texttt{quant::pde::thomas} in
C++). It will be reused by Crank--Nicolson in Block~3 and by PSOR in
Block~4.

% =====================================================================
\section{The BTCS discretisation}
% =====================================================================

We discretise the transformed Black--Scholes PDE
\begin{equation}
\frac{\partial V}{\partial \tau}
= \frac{\sigma^2}{2}\frac{\partial^2 V}{\partial x^2}
+ \mu\,\frac{\partial V}{\partial x}
- r\,V
\label{eq:transformed}
\end{equation}
on the grid of Block~0. BTCS evaluates every spatial stencil at the
\emph{new} time level $n+1$:
\[
\frac{V_j^{n+1} - V_j^n}{\Delta\tau}
= \frac{\sigma^2}{2}\,\frac{V_{j+1}^{n+1} - 2 V_j^{n+1} + V_{j-1}^{n+1}}{\Delta x^2}
+ \mu\,\frac{V_{j+1}^{n+1} - V_{j-1}^{n+1}}{2\Delta x}
- r\,V_j^{n+1}.
\]
Multiplying by $\Delta\tau$ and grouping all $V^{n+1}$ on the left,
$V^n$ on the right:
\begin{equation}
\boxed{\;
b_-\,V_{j-1}^{n+1} + b_0\,V_j^{n+1} + b_+\,V_{j+1}^{n+1}
\;=\; V_j^n
\;}
\label{eq:btcs}
\end{equation}
with
\[
b_- = -\alpha + \frac{\mu\,\Delta\tau}{2\,\Delta x},
\qquad
b_0 = 1 + 2\alpha + r\,\Delta\tau,
\qquad
b_+ = -\alpha - \frac{\mu\,\Delta\tau}{2\,\Delta x},
\]
where $\alpha = (\sigma^2/2)\Delta\tau/\Delta x^2$ as before.

Equation \eqref{eq:btcs} is the FTCS scheme \emph{written from the
other side}: where FTCS reads ``$V_j^{n+1}$ equals an explicit
combination of $V^n$'', BTCS reads ``$V_j^n$ equals an explicit
combination of $V^{n+1}$''. The latter is a system to be solved for
the unknowns $V^{n+1}$. The two schemes are the $\theta = 0$ and
$\theta = 1$ extremes of the family
\[
b_-\,V_{j-1}^{n+1} + b_0\,V_j^{n+1} + b_+\,V_{j+1}^{n+1}
= a_-\,V_{j-1}^n + a_0\,V_j^n + a_+\,V_{j+1}^n,
\]
the $\theta$-scheme, with Crank--Nicolson at $\theta = 1/2$ in
Block~3.

% =====================================================================
\section{The tridiagonal system}
% =====================================================================

For each time step, equation \eqref{eq:btcs} written for the interior
nodes $j = 1, \dots, N-1$ is a linear system $A\,\mathbf{V}^{n+1}_{\mathrm{int}} = \mathbf{d}$
where $A \in \mathbb{R}^{(N-1)\times(N-1)}$ is tridiagonal:
\[
A = \begin{pmatrix}
b_0 & b_+ &        &        & \\
b_- & b_0 & b_+    &        & \\
    & b_- & b_0    & b_+    & \\
    &     & \ddots & \ddots & \ddots \\
    &     &        & b_-    & b_0
\end{pmatrix}.
\]
The boundary nodes $V_0^{n+1}$ and $V_N^{n+1}$ are known from the
Dirichlet conditions of Block~0; they appear in the right-hand side:
\begin{equation}
\mathbf{d} = \begin{pmatrix}
V_1^n - b_-\,V_0^{n+1} \\
V_2^n \\
\vdots \\
V_{N-2}^n \\
V_{N-1}^n - b_+\,V_N^{n+1}
\end{pmatrix}.
\label{eq:rhs}
\end{equation}
The first and last rows pick up the contributions from the
boundary; the rest of $\mathbf{d}$ is just $V^n$ at the interior.

\paragraph{Constancy of $A$ in time.} The coefficients $b_-, b_0, b_+$
are independent of $n$ and $j$ thanks to the constant-coefficient
form \eqref{eq:transformed}. The matrix $A$ is therefore the same at
every time step; only the right-hand side $\mathbf{d}$ changes.
This is the structural pay-off of the logarithmic transformation
that we already exploited in Block~1, now in a stronger form: at
each step we solve a system with the \emph{same} matrix and a
different right-hand side. The forward sweep of the Thomas algorithm
can be done \emph{once} at setup and reused, halving the per-step
cost.

% =====================================================================
\section{Stability: unconditional}
% =====================================================================

The von Neumann analysis of \eqref{eq:btcs} mirrors Block 0, §7.3
extended with the convective and reactive terms. Substituting
$V_j^n = g^n e^{ij\theta}$ and solving for $g$:
\begin{equation}
g(\theta) = \frac{1}{1 + 4\alpha\sin^2(\theta/2) - i\,\beta\sin\theta + r\,\Delta\tau},
\qquad \beta = \frac{\mu\,\Delta\tau}{\Delta x}.
\label{eq:btcs-amplification}
\end{equation}
The squared modulus,
\[
|g|^2 = \frac{1}{\bigl(1 + 4\alpha\sin^2(\theta/2) + r\Delta\tau\bigr)^2 + \beta^2\sin^2\theta},
\]
has denominator $\ge 1$ for every $\theta$ (every term is
non-negative; one of them, $1$, is strictly so). Hence $|g(\theta)|
\le 1$ unconditionally: BTCS is stable for any $\Delta\tau > 0$ and
$\Delta x > 0$.

Combined with consistency ($O(\Delta\tau) + O(\Delta x^2)$ local
truncation), the Lax equivalence theorem gives convergence:
\[
\|V_h - V\|_\infty = O(\Delta\tau) + O(\Delta x^2).
\]

\begin{remark}[Stability is not monotonicity]
$L^2$ stability does not imply $L^\infty$ monotonicity. For very
large $\Delta\tau$ at moderate $\Delta x$, the matrix $A$ ceases to
be an M-matrix: $A^{-1}$ has negative entries, and the discrete
operator can produce localised oscillations near the strike kink.
These oscillations are bounded (stability holds), but they degrade
local accuracy. For our parameter range this regime is far away;
when we reach it (large-$\Delta\tau$ tests in the validation), we
will see prices that are merely imprecise, never explosive.
\end{remark}

% =====================================================================
\section{The Thomas algorithm}
% =====================================================================

Solving $A\,\mathbf{x} = \mathbf{d}$ with general Gaussian elimination
costs $O(N^3)$ — prohibitive at $N = 200$, hopeless at $N = 1000$.
The Thomas algorithm exploits the tridiagonal structure to compute
the same solution in $O(N)$ work, which is asymptotically optimal:
just reading the input takes $O(N)$ time.

\paragraph{Setup.} Consider the general tridiagonal system
\[
\begin{pmatrix}
\tilde b_1 & c_1        &            &            &        \\
a_2        & \tilde b_2 & c_2        &            &        \\
           & a_3        & \tilde b_3 & c_3        &        \\
           &            & \ddots     & \ddots     & \ddots \\
           &            &            & a_n        & \tilde b_n
\end{pmatrix}
\begin{pmatrix} x_1 \\ x_2 \\ \vdots \\ x_n \end{pmatrix}
=
\begin{pmatrix} d_1 \\ d_2 \\ \vdots \\ d_n \end{pmatrix}.
\]
We write $\tilde b$ for the diagonal of the general system to
distinguish from the BTCS-specific $b_0, b_-, b_+$.

\paragraph{Forward sweep.} Standard Gaussian elimination, applied
naively to a tridiagonal system, produces \emph{no fill}: at each
step the only non-zero entries below the pivot are $a_i$, and the
elimination touches $\tilde b_i$ and $d_i$ but not the rows below.
The recursion is
\begin{align}
c'_1 &= \frac{c_1}{\tilde b_1}, &
d'_1 &= \frac{d_1}{\tilde b_1}, \nonumber \\
m_i  &= \tilde b_i - a_i\,c'_{i-1}, &
c'_i &= \frac{c_i}{m_i}, &
d'_i &= \frac{d_i - a_i\,d'_{i-1}}{m_i}, \quad i = 2, \dots, n.
\label{eq:thomas-forward}
\end{align}
The $m_i$ are the \emph{effective pivots}: what remains on the
diagonal after eliminating $a_i$. After the sweep the system is
upper bidiagonal:
\[
\begin{pmatrix}
1 & c'_1 &       &       \\
  & 1    & c'_2  &       \\
  &      & 1     & c'_3  \\
  &      &       & \ddots \\
\end{pmatrix}
\mathbf{x} = \mathbf{d'}.
\]

\paragraph{Backward substitution.} Trivial:
\begin{equation}
x_n = d'_n, \qquad
x_i = d'_i - c'_i\,x_{i+1}, \quad i = n-1, n-2, \dots, 1.
\label{eq:thomas-back}
\end{equation}

\paragraph{Operation count.} The forward sweep does roughly $5n$
floating-point operations (one division for $c'_i$, three operations
for $d'_i$, one for $m_i$). The backward substitution does $\sim 2n$.
Total: $O(n)$ with constant $\approx 7$.

\paragraph{Memory.} Only four vectors of length $\le n$: the three
diagonals (modified in place if desired) and the right-hand side.
No fill-in, no $O(n^2)$ storage.

\paragraph{Pre-factoring for repeated solves.} The forward sweep
\eqref{eq:thomas-forward} computes two distinct families of
quantities: those depending only on $(a, \tilde b, c)$, namely
$\{m_i, c'_i\}$, and those depending also on $d$, namely
$\{d'_i\}$. When the matrix is constant across many right-hand
sides — exactly our case, with one fixed $A$ and $M$ different
right-hand sides — we compute the $(m_i, c'_i)$ once at setup, and
each subsequent solve costs only the $d'$ recursion plus the
backward substitution. The savings are $O(n)$ divisions per solve,
i.e., a factor $\sim 1.5$--$2$ in practice.

% =====================================================================
\section{Stability of Thomas without pivoting}
% =====================================================================

The Thomas algorithm as written does \emph{not} pivot: there is no
provision for $m_i = 0$ or $|m_i|$ small. In general, Gaussian
elimination without pivoting is unsafe — a small pivot amplifies
round-off catastrophically. For Thomas there is a clean theorem
that rules out the failure mode for our matrices.

\begin{theorem}[Thomas without pivoting under diagonal dominance]
\label{thm:thomas-stab}
If the tridiagonal matrix is diagonally dominant in the sense
\begin{equation}
|\tilde b_i| \ge |a_i| + |c_i|, \quad i = 1, \dots, n,
\label{eq:diag-dominance}
\end{equation}
with strict inequality in at least one row, then all effective
pivots $m_i$ in \eqref{eq:thomas-forward} satisfy $|m_i| \ge
|\tilde b_i| - |a_i| > 0$, and the algorithm is backward stable:
the computed solution $\hat{\mathbf{x}}$ satisfies
$(A + \delta A)\hat{\mathbf{x}} = \mathbf{d}$ with
$\|\delta A\|_\infty \le c_n\,\varepsilon_{\mathrm{mach}}\|A\|_\infty$
for a small constant $c_n$ proportional to $n$.
\end{theorem}

(Proof in Higham, \emph{Accuracy and Stability of Numerical Algorithms},
\S9.5.)

\paragraph{Application to BTCS.} The BTCS matrix has $\tilde b_i =
b_0$, $a_i = b_-$, $c_i = b_+$, all constant in $i$. Diagonal
dominance reduces to
\[
|b_0| \;\ge\; |b_-| + |b_+|.
\]
We have $|b_0| = 1 + 2\alpha + r\Delta\tau$ (positive). Computing the
right-hand side:
\[
|b_-| + |b_+|
= \left|\alpha - \frac{\mu\Delta\tau}{2\Delta x}\right|
 + \left|\alpha + \frac{\mu\Delta\tau}{2\Delta x}\right|.
\]
When $\alpha \ge \bigl|\mu\Delta\tau/(2\Delta x)\bigr|$ — i.e.,
$\bigl|\mu\bigr|\Delta x/\sigma^2 \le 1$, the grid Péclet condition
of Block~1 — both terms inside the absolute values have $\alpha$
dominating and the same sign, and we get
$|b_-| + |b_+| = 2\alpha$. Then
\[
|b_0| - (|b_-| + |b_+|) = 1 + r\Delta\tau > 0,
\]
strict diagonal dominance. So Thomas without pivoting is provably
stable for BTCS whenever $\mathrm{Pe}_h \le 1$, which the project's
parameter range satisfies with comfortable margin
($\mathrm{Pe}_h \le 0.25$).

% =====================================================================
\section{Convergence rates and computational cost}
% =====================================================================

The BTCS error is $O(\Delta\tau) + O(\Delta x^2)$. With $M$ now
independent of $N$, the user has a choice:

\paragraph{Quadratic-precision regime ($M \sim N^2$).} Choose
$\Delta\tau \sim \Delta x^2$. Both error contributions are
$O(\Delta x^2)$; the total error is $O(\Delta x^2)$. The
computational cost is $O(N) \cdot M \cdot 2 = O(N \cdot N^2) =
O(N^3)$, with the constant a factor $\sim 7$ from the Thomas count.
\emph{Same asymptotic cost as FTCS for the same precision.}

\paragraph{Linear-precision regime ($M \sim N$).} Choose
$\Delta\tau \sim \Delta x$. The error is dominated by $O(\Delta\tau)
= O(\Delta x)$; the spatial $O(\Delta x^2)$ contribution is
sub-dominant. The total error is $O(\Delta x)$. Cost: $O(N) \cdot
M = O(N^2)$. \emph{Acceptable when first-order precision is enough}
(prototyping, calibration of insensitive parameters, sanity checks).

\paragraph{Crank--Nicolson preview.} The next block achieves
$O(\Delta\tau^2) + O(\Delta x^2)$ local truncation, so $M \sim N$
gives full $O(\Delta x^2)$ precision at $O(N^2)$ cost. That is the
point where the explicit-vs-implicit trade decisively tips toward
implicit.

\paragraph{Summary table.}
\begin{center}
\begin{tabular}{lcccc}
\hline
Scheme & Local error & Stability & For $O(\Delta x^2)$ & For $O(\Delta x)$ \\
\hline
FTCS  & $\Delta\tau + \Delta x^2$  & $\alpha \le 1/2$       & $O(N^3)$ (forced) & not relevant \\
BTCS  & $\Delta\tau + \Delta x^2$  & unconditional          & $O(N^3)$ (M $\sim N^2$) & $O(N^2)$ (M $\sim N$) \\
CN    & $\Delta\tau^2 + \Delta x^2$ & unconditional         & $O(N^2)$ (M $\sim N$) & overkill \\
\hline
\end{tabular}
\end{center}

% =====================================================================
\section{Implementation outline}
% =====================================================================

The deliverables are two modules each in Python and C++:

\begin{itemize}[leftmargin=2em]
\item \texttt{quantlib.thomas} (Python) / \texttt{quant::pde::thomas\_*}
(C++) provides
\texttt{thomas\_solve(sub, diag, sup, rhs)} for one-shot solves and
the pair \texttt{thomas\_factor(sub, diag, sup)} +
\texttt{thomas\_solve\_factored(factor, rhs)} for repeated solves
with a fixed matrix.
\item \texttt{quantlib.btcs} (Python) / \texttt{quant::pde::btcs\_*}
(C++) provides \texttt{btcs\_european\_call/put} with the same
signature as the FTCS pricers, except that there is no CFL check:
any $M \ge 1$ is valid input.
\end{itemize}

The validation suite checks five properties:
\begin{enumerate}[leftmargin=2em]
\item Cross-validation of \texttt{thomas\_solve} against
\texttt{numpy.linalg.solve} (Python) and against an independent dense
LU implementation in the C++ tests, for randomly-generated diagonally
dominant matrices.
\item Cross-validation of BTCS against Black--Scholes closed-form,
with tolerance scaled to the discretisation error.
\item Convergence rate at $\Delta\tau \sim \Delta x^2$: error reduction
by a factor of $\sim 4$ when $N \to 2N$, $M \to 4M$.
\item First-order convergence in time at fixed (large) $N$, $M \to
2M$: error reduction by $\sim 2$.
\item \emph{Large-$\Delta\tau$ stability test}: BTCS at $M = 10$
for $T = 1$ produces a finite, sensible price (although imprecise),
with no explosion. This is the operational signature of unconditional
stability.
\end{enumerate}

% =====================================================================
\section*{References}
% =====================================================================

\begin{thebibliography}{9}

\bibitem{TavellaRandall2000}
D. Tavella and C. Randall.
\emph{Pricing Financial Instruments: The Finite Difference Method.}
Wiley, 2000. \S3.2 (BTCS for BS), \S3.4 (stability of implicit
schemes).

\bibitem{Seydel2017}
R. Seydel.
\emph{Tools for Computational Finance.}
6th ed., Springer, 2017. \S4.2.2 (BTCS), \S4.6 (linear solvers).

\bibitem{Strikwerda2004}
J.~C. Strikwerda.
\emph{Finite Difference Schemes and Partial Differential Equations.}
2nd ed., SIAM, 2004. \S3.1--3.2 (implicit methods, von Neumann
analysis).

\bibitem{NumericalRecipes}
W.~H. Press, S.~A. Teukolsky, W.~T. Vetterling and B.~P. Flannery.
\emph{Numerical Recipes: The Art of Scientific Computing.}
3rd ed., CUP, 2007. \S2.4 (the Thomas algorithm with code).

\bibitem{Higham2002}
N.~J. Higham.
\emph{Accuracy and Stability of Numerical Algorithms.}
2nd ed., SIAM, 2002. \S9.5 (stability of LU without pivoting for
diagonally-dominant matrices, including the proof of
Theorem~\ref{thm:thomas-stab}).

\bibitem{TrefethenBau1997}
L.~N. Trefethen and D.~Bau~III.
\emph{Numerical Linear Algebra.}
SIAM, 1997. Lecture~21 (Gaussian elimination without pivoting).

\bibitem{Smith1985}
G.~D. Smith.
\emph{Numerical Solution of Partial Differential Equations:
Finite Difference Methods.}
3rd ed., OUP, 1985. \S1.13 (the $\theta$-family of schemes).

\end{thebibliography}

\end{document}
